
\documentclass{article}
\usepackage{colortbl}
\usepackage{makecell}
\usepackage{multirow}
\usepackage{supertabular}

\begin{document}

\newcounter{utterance}

\centering \large Interaction Transcript for game `cladder', experiment `full\_v1.5\_default', episode 9292 with llama{-}3.1{-}8b{-}instruct.
\vspace{24pt}

{ \footnotesize  \setcounter{utterance}{1}
\setlength{\tabcolsep}{0pt}
\begin{supertabular}{c@{$\;$}|p{.15\linewidth}@{}p{.15\linewidth}p{.15\linewidth}p{.15\linewidth}p{.15\linewidth}p{.15\linewidth}}
    \# & \multicolumn{2}{c}{Player} && \multicolumn{2}{c}{Game Master} \\
    \hline

    \theutterance \stepcounter{utterance}  
    & & & \multicolumn{4}{p{0.6\linewidth}}{
        \cellcolor[rgb]{0.9,0.9,0.9}{
            \makecell[{{p{\linewidth}}}]{
                \texttt{\tiny{[P1$\langle$GM]}}
                \texttt{You are an expert in causal inference. The following question is not a typical commonsense query, but rather a meticulously designed question created by a professor specializing in causal inference, intended to assess the students' mastery of the course content.} \\
\texttt{Imagine a self{-}contained, hypothetical world with only the following conditions, and without any unmentioned factors or causal relationships: Zory has a direct effect on zuph and glimx. Zuph has a direct effect on jyka. Jyka has a direct effect on glimx. Zory is unobserved. The overall probability of zuph is 36\%. The probability of not zuph and glimx is 30\%. The probability of zuph and glimx is 31\%. Is the chance of glimx smaller when observing zuph?} \\
\texttt{Start your answer with 'Yes' or 'No', followed by additional reasoning or evidence to support your explanation.} \\
            }
        }
    }
     \\ \\

    \theutterance \stepcounter{utterance}  
    & \multicolumn{4}{p{0.6\linewidth}}{
        \cellcolor[rgb]{0.561,0.737,0.561}{
            \makecell[{{p{\linewidth}}}]{
                \texttt{\tiny{[P1$\rangle$GM]}}
                \texttt{No.} \\
\\ 
\texttt{To determine whether the chance of glimx is smaller when observing zuph, we need to consider the relationships between the variables and the given probabilities.} \\
\\ 
\texttt{Given:} \\
\texttt{{-} Zory has a direct effect on zuph and glimx.} \\
\texttt{{-} Zuph has a direct effect on jyka.} \\
\texttt{{-} Jyka has a direct effect on glimx.} \\
\texttt{{-} Zory is unobserved.} \\
\texttt{{-} P(zuph) = 36\%} \\
\texttt{{-} P(not zuph and glimx) = 30\%} \\
\texttt{{-} P(zuph and glimx) = 31\%} \\
\\ 
\texttt{We can use the law of total probability to find the probability of glimx. Let's denote the probability of glimx as P(glimx).} \\
\\ 
\texttt{P(glimx) = P(glimx | zuph) * P(zuph) + P(glimx | not zuph)} \\
\\ 
\texttt{We know that P(zuph) = 36\%. To find P(glimx | zuph), we can use the fact that P(zuph and glimx) = 31\% and P(zuph) = 36\%. Since P(zuph and glimx) is a subset of P(zuph), we can write:} \\
\\ 
\texttt{P(glimx | zuph) = P(zuph and glimx) / P(zuph) = 31\% / 36\% ≈ 0.8611} \\
\\ 
\texttt{Now, we need to find P(glimx | not zuph). We know that P(not zuph and glimx) = 30\%. Since P(not zuph and glimx) is a subset of P(not zuph), we can write:} \\
\\ 
\texttt{P(glimx | not zuph) = P(not zuph and glimx) / P(not zuph)} \\
\\ 
\texttt{To find P(not zuph), we can use the fact that P(zuph) = 36\% and P(zuph or not zuph) = 1. Therefore, P(not zuph) = 1 {-} P(zuph) = 1 {-} 36\% = 64\%.} \\
\\ 
\texttt{P(glimx | not zuph) = P(not zuph and glimx) / P(not zuph) = 30\% / 64\% ≈ 0.4688} \\
\\ 
\texttt{Now, we can find P(glimx) using the law of total probability:} \\
\\ 
\texttt{P(glimx) = P(glimx | zuph) * P(zuph) + P(glimx | not zuph) * P(not zuph)} \\
\texttt{= 0.8611 * 0.36 + 0.4688 * 0.64} \\
\texttt{= 0.3096 + 0.3003} \\
\texttt{= 0.6099} \\
\\ 
\texttt{To determine whether} \\
            }
        }
    }
    & & \\ \\

    \theutterance \stepcounter{utterance}  
    & & & \multicolumn{2}{p{0.3\linewidth}}{
        \cellcolor[rgb]{0.9,0.9,0.9}{
            \makecell[{{p{\linewidth}}}]{
                \texttt{\tiny{[GM$|$GM]}}
                \texttt{No} \\
            }
        }
    }
    & & \\ \\

    \theutterance \stepcounter{utterance}  
    & & & \multicolumn{2}{p{0.3\linewidth}}{
        \cellcolor[rgb]{0.9,0.9,0.9}{
            \makecell[{{p{\linewidth}}}]{
                \texttt{\tiny{[GM$|$GM]}}
                \texttt{no} \\
            }
        }
    }
    & & \\ \\

    \theutterance \stepcounter{utterance}  
    & & & \multicolumn{2}{p{0.3\linewidth}}{
        \cellcolor[rgb]{0.9,0.9,0.9}{
            \makecell[{{p{\linewidth}}}]{
                \texttt{\tiny{[GM$|$GM]}}
                \texttt{game\_result = WIN} \\
            }
        }
    }
    & & \\ \\

\end{supertabular}
}

\end{document}
